\documentclass[12 pt, letterpaper]{hmcpset}
\usepackage{hw-preamble}

\name{Jayden Thadani}
\class{MATH 131 --- Mathematical Analysis I}
\assignment{Homework 10}
\duedate{23rd November 2024}

\begin{document}

\begin{problem}[Rudin chapter 5, exercise 9]
	Let $f$ be a continuous real function on $\mathbb{R}$, of which it is known that $f'(x)$ exists for all $x \neq 0$ and that $f'(x) \to 3$ as $x \to 0$. Does it follow that $f'(0)$ exists?
\end{problem}

\begin{solution}
	We know that as $x \to 0$, $f'(x) \to 3$. Consider
	\[
		\lim_{h \to 0} \frac{f(x + h) - f(x)}{h}.
	\]
	Since $f$ is continuous, the numerator and denominator approach $0$ as $h \to 0$. Thus, we can apply L'H\^opital's rule (differentiating with respect to $h$) to get
	\begin{equation*}
		\lim_{h \to 0} \frac{f(0 + h) - f(0)}{h} = \lim_{h \to 0} \frac{f'(0 + h) - 0}{1} = \lim_{h \to 0} f'(h) = 3.
	\end{equation*}
	as desired.
\end{solution}

\pagebreak

\begin{problem}[Rudin chapter 5, exercise 11]
	Suppose $f$ is defined in a neighbourhood of $x$ and suppose $f''(x)$ exists. Show that
	\[
		\lim_{h \to 0} \frac{f(x + h) + f(x - h) - 2 f(x)}{h^{2}} = f''(x).
	\]
	Show by an example that the limit may exist even if $f''(x)$ does not.
\end{problem}

\begin{solution}
	By L'H\^opital's rule we can differentiate with respect to $h$ above and below. That is,
	\begin{align*}
		\lim_{h \to 0} \frac{f(x + h) + f(x - h) - 2 f(x)}{h^{2}} & = \lim_{h \to 0} \frac{f'(x + h) - f(x - h)}{2 h} \\
		& = \lim_{h \to 0} \frac{f'(x + h) - f(x)}{2h} + \lim_{-h \to 0} \frac{f'(x) - f(x + h)}{-2h} \\
		& = \frac{f''(x)}{2} + \lim_{h \to 0} \frac{f'(x + h) - f'(x)}{2h} \\
		& = \frac{f''(x)}{2} + \frac{f''(x)}{2} \\
		& = f''(x).
	\end{align*}
	
	Consider $f : \mathbb{R} \to \mathbb{R}$ by $x \mapsto x / \lvert x \rvert$ and $0 \mapsto 0$. Then since the numerator of the quotient $f(0 + h) + f(0 - h) - 2 f(0) = 0$ for all $h \in \mathbb{R}$, the limit is
	\[
		\lim_{h \to 0} \frac{f(0 + h) + f(0 - h) - 2 f(0)}{h^{2}} = 0
	\]
	However, $f$ is not even once differentiable, much less twice differentiable at $0$ --- the left and right limits of the quotient
	\[
		\frac{f(x + h) - f(x)}{h}
	\]
	are different for $x > 0$ and $x < 0$, and thus, the derivative does not exist.
\end{solution}

\pagebreak

\begin{problem}[Rudin chapter 5, exercise 12]
	If $f(x) = \lvert x \rvert^{3}$, compute $f'(x)$, $f''(x)$ for all real $x$, and show that $f^{(3)}(0)$ does not exist.	
\end{problem}

\begin{solution}
	Note that by exercise 9 (the first exercise on this homework) it suffices just to calculate $f'$ and $f''$ for all $x \neq 0$ and show that there is a defined limit as $x \to 0$.

	For $x > 0$, $f(x) = x^{3}$, and thus, $f'(x) = 3 x^{2}$. For $x < 0$, $f(x) = - x^{3}$, and thus, $f'(x) = -3 x^{2}$. Since $f'(x) \to 0$ as $x \to 0$, we have $f$ differentiable with $f'(0) = 0$. That is,
	\[
		f'(x) = \begin{cases}
			- 3 x^{2} & x < 0 \\
			3 x^{2} & x \ge 0.
		\end{cases}
	\]

	Thus,  $f''(x) = - 6 x$ for $x < 0$ and $f''(x) = 6 x$ for $x > 0$. Further, since $f''(x) \to 0$ if $x \to 0$. That is,
	\[
		f''(x) = 6 \lvert x \rvert.
	\]

	However this has derivative $f^{(3)}(x) = -6$ for $x < 0$ and $f^{(3)}(x) = 6$ for $x > 0$. Since a differentiable function cannot have a derivative with a jump discontinuity, we must conclude that $f^{(3)}(0)$ does not exist.
\end{solution}

\pagebreak

\begin{problem}[Rudin chapter 5, exercise 13]
	Suppose $a$ and $c$ are real numbers $c > 0$, and $f$ is defined on $[-1, 1]$ by
	\[
		f(x) = \begin{cases}
			x^{a} \sin (\lvert x \rvert^{-c}) & \text{if } x \neq 0 \\
			0 & \text{if } x = 0.
		\end{cases}
	\]
	Prove the following statements
	\begin{enumerate}
		\item $f$ is continuous if and only if $a > 0$.
		\item $f'(0)$ exists if and only if $a > 1$.
		\item $f'$ is bounded if and only if $a \ge 1 + c$.
		\item $f'$ is continuous if and only if $a > 1 + c$.
		\item $f''(0)$ exists if and only if $a > 2 + c$.
		\item $f''$ is bounded if and only if $a \ge 2 + 2c$.
		\item $f''$ is continuous if and only if $a > 2 + 2c$.
	\end{enumerate}
\end{problem}

\begin{solution}
	\begin{enumerate}
		\item Suppose $a > 0$. Then $\lim_{x \to 0} x^{a} = 0$. But since $\sin{\lvert x \rvert^{-c}}$ is bounded within $[-1, 1]$, this means that as $x \to 0$, $f(x)$ gets arbitrarily small --- $f(x) \to 0$. That is, $\lim_{x \to 0} f(x) = f(0)$. Since $f$ is the composition of functions that are continuous everywhere except $0$, this is sufficient to establish that $f$ is continuous on $\mathbb{R}$.

			If $a = 0$, then it is well known that $f(x) = \sin{(\lvert x \rvert^{-c})}$ is discontinuous. This follows since in any small neighbourhood $(- \delta, \delta)$ around $0$ there is some $x \in (-\delta, \delta)$ with $1/\lvert x \rvert^{c} = (2m + 1) \pi / 2 > 1 / \delta$ for sufficiently large $m$, and thus $\sin{(1/\lvert x \rvert^{c})} = 1$. Thus, $f(x) \not \to 0$ as $x \to 0$. That is, $f$ is not continuous.

			If $a < 0$, then $x^{a}$ is unbounded in every neighbourhood of $0$ (by the Archimedean property, as $n^{-a} \to \infty$ we also have $n^{a} \to 0$ and the result follows by sequential limits). Since $\sin{(1 / \lvert x \rvert^{c})} \not \to 0$, the product $f(x)$ is also unbounded near $0$. Thus, $f$ cannot be continuous.
		\item For $x \neq 0$, $f$ is differentiable with
			\[
				f'(x) = \begin{cases}
					a x^{a - 1} \sin{(\lvert x \rvert^{-c})} - c x^{a - c - 1} \cos{(x^{-c})} & x < 0 \\
					a x^{a - 1} \sin{(\lvert x \rvert^{-c})} + c x^{a - c - 1} \cos{(-x^{-c})} & x > 0
				\end{cases}
			\]
			By similar analysis as in part (a), when $a > 1$ the limits agree.
		\item Using the expression for $f'$ in the previous part, the $x^{a - c - 1}$ is only bounded when $a \ge 1 + c$, so the derivative is only bounded when $a \ge 1 + c$ --- then the function is just bounded by the fact that $\sin{}$ and $\cos{}$ are bounded.
		\item For $x \neq 0$, $f'$ is differentiable with
			\[
				f''(x) = a x^{a - 2} \sin{(\lvert x^{-c} \rvert)} + a x^{a - 1} \cos{\lvert x \rvert^{-c}} + c x^{a - c - 2} \cos{(x^{-c})} 
			\]
		$f''(0)$ only exists when the $x^{a - 2 - c}$ term has positive exponent in front of the $\sin$ function --- exactly when $a > 2 + c$.
	\item For boundedness, we consider the $\cos{}$ term which needs $a \ge 2 + 2c$ for the $x^{a - 2 - 2c}$ term in front of it to have non-negative exponent after which the function is again bounded by the boundedness of trigonometric functions.
	\item Again, if the $x^{a - 2 - 2c}$ term has strictly positive exponenent, all trigonometric functions are dominated by the polynomial terms and are thus continuous at $x = 0$.
	\end{enumerate}
\end{solution}

\pagebreak

\begin{problem}[Rudin chapter 5, exercise 14]
	Let $f$ be a differentiable real function defined in $(a, b)$. Prove that $f$ is convex if and only if $f'$ is monotonically increasing. Assume next that $f''(x)$ exists for every $x \in (a, b)$ and prove that $f$ is convex if and only if $f''(x) \ge 0$ for all $x \in (a, b)$.
\end{problem}

\begin{solution}
	Suppose $f'$ is monotonically increasing. Suppose $x, y \in (a, b)$ and without loss of generality, assume $x < y$. (If $x > y$, note that the property still holds by replacing $\lambda$ by $1 - \lambda$). Now, for any $z = \lambda x + (1 - \lambda) y$, we want to show $f(z) \le \lambda f(x) + (1 - \lambda) f(y)$. Equivalently, we want to show $f(z) - f(y) \le \lambda (f(x) - f(y))$.

	Note that $\lambda(x - y) + y = z$ so $\lambda = (z - y) / (x - y)$. But now (with some algebraic rearrangement) the inequality we want to show is just
	\[
		\frac{f(y) - f(z)}{y - z} \ge \frac{f(z) - f(x)}{z - x}
	\]
	which is true by applying the mean value theorem and the fact that $f'$ is monotonically increasing.

	We have shown that a differentiable $f$ is convex if and only if $f'$ is monotonically increasing. If $f$ is twice differentiable, then $f'$ is monotonically increasing if and only if $f''(x) \ge 0$ for all $x \in (a, b)$.
\end{solution}

\end{document}
